% ============================================================
%  THE ORTYX PROTOCOL
%  Part IV — Reconstruction
%  Chapter 17: Signal, Salience, and Constraint
% ============================================================

\chapter{Signal, Salience, and Constraint}
\label{ch:signal-salience-constraint}

\chapepigraph{What the auditory system selects from a continuous
stream is not the loudest, not the newest, not the most complex.
It is what fits the smallest admissible trajectory space:
what is most constrained, most predictable, most organized.
Whether this is what \emph{matters} is the question the
framework cannot answer for us.}{Flyxion}

\noindent
The Marine Salience Algorithm, stripped of the
consciousness claims and the authenticity ontology,
is a well-engineered time-domain coherence detector
with real applications in audio quality assessment
and perceptual computing. Its translation into
RSVP-compatible terms does not change this. What
the translation does is make visible the theoretical
assumptions that the original formulation carried
implicitly.

\section{Jitter as Trajectory Constraint Density}
\label{sec:jitter-trajectory}

In the original Phoenix formulation, period jitter
$\Jp(i) = |T_i - \bar{T}_i|$ is described as
measuring the ``incoherence'' of a signal at peak $i$.
This description is accurate but understates the
geometric content of the metric. In RSVP-compatible
terms, jitter measures the local constraint density
of the signal's trajectory in its state space.

Consider the signal $s(t)$ as tracing a trajectory
through a state space $X$. At each moment, the
admissible future trajectories from the current state
form a set $\AdmTraj(x_t) \subseteq X$. The
accessibility entropy is:
\[
  \AccEnt(x_t) = \log |\AdmTraj(x_t)|.
\]
A signal with low jitter --- one whose inter-peak
intervals are stable and predictable --- is a signal
whose current state highly constrains future states:
the admissible trajectory space is small. A signal
with high jitter has many admissible continuations;
its current state constrains future states weakly.

The Marine metric thus approximates the inverse
of accessibility entropy:
\[
  \SalFunc(i) \propto \frac{1}{\AccEnt(x_{t_i})}.
\]
This translation makes explicit what the original
formulation implied: the Marine is detecting signals
that occupy low-entropy, highly constrained regions
of trajectory space. These are regions where the
signal's dynamics are, in the information-theoretic
sense, most determined by its current state.

\begin{technote}[Technical Note: Accessibility and Jitter]
The proportionality $\SalFunc \propto 1/\AccEnt$
is a structural observation rather than a derivation.
To make it a derivation, one would need to specify:
(1)~the state space $X$ for audio signals,
(2)~a measure on $\AdmTraj(x_t)$ that yields a
well-defined accessibility entropy, and
(3)~a demonstration that low jitter signals, as
detected by the Marine EMA procedure, correspond
to signals in low-$\AccEnt$ regions. Steps~(1)
and~(2) require theoretical work not yet completed.
Step~(3) requires empirical verification. The
translation is therefore at \LevelII{} (productive
metaphor generating a research direction) rather
than \LevelI{} (established engineering result).
\end{technote}

\section{The Salience Filter as Constraint-Governed Gating}
\label{sec:salience-gating}

The Marine filter's selection of high-salience events
for downstream processing becomes, in RSVP-compatible
terms, a constraint-governed gating mechanism. The
filter passes events that are consistent with the
signal's running admissibility model: events that
the model predicted were possible, that arrive
within the model's tolerance for timing and amplitude,
and that therefore require minimal trajectory revision
to incorporate.

This translation reveals the connection to predictive
processing that was noted as a structural observation
in Chapter~5 and classified as a reconstructive
conjecture. In the RSVP translation, the connection
becomes more precise: both the Marine filter and the
predictive processing framework are selecting events
based on their compatibility with a running model
of the signal's admissible trajectory space. The
Marine filter does this with the specific EMA-based
procedure; predictive processing accounts do it with
hierarchical generative models. The shared structure
is constraint-governed selection.

\begin{salvage}[Reconstruction: Salience as Admissibility Test]
A constraint-preserving reformulation of the Marine
salience filter: given a signal $s(t)$ and a running
model $\mathcal{M}_t$ of the signal's admissible
trajectory space, the salience of event $e_i$ at
time $t_i$ is:
\[
  \SalFunc(i)
  = \mathbf{1}\bigl(e_i \in \AdmTraj(\mathcal{M}_{t_i-1})\bigr)
  \cdot \left(1 - \frac{\AccEnt(\mathcal{M}_{t_i})}
    {\AccEnt(\mathcal{M}_{t_i-1})}\right).
\]
High-salience events are those consistent with the
running model ($\mathbf{1} = 1$) that reduce the
accessibility entropy of the model (the fraction
term is large). This formulation makes explicit
what the Marine EMA procedure approximates: the
detection of events that are both predicted and
constraining.
\end{salvage}

\section{Harmonic Conformity as Constraint Propagation}
\label{sec:conformity-propagation}

The \HCFone{} conformity relation $\Gamma_n(t) =
\mathbf{1}(|v_n(t) - v_0(t)| < \varepsilon)$ acquires
a natural RSVP interpretation as a constraint propagation
test. Two harmonics $n$ and $0$ satisfy the conformity
relation at time $t$ when they are evolving within
a shared admissible trajectory basin: the constraints
governing the fundamental's evolution are also governing
the harmonic's evolution within tolerance $\varepsilon$.

This interpretation connects the \HCFone{} conformity
relation to the broader RSVP principle that physically
coupled systems --- systems whose dynamics are governed
by shared constraints --- can be represented efficiently
by encoding the constraints once rather than the
trajectories separately. The $O(1) + O(k)$ encoding
complexity of consensus encoding is a consequence
of constraint sharing: most harmonics are in the same
constraint basin as the fundamental, requiring only
a binary conformity flag rather than independent
trajectory specification.

\begin{auditprop}
\label{ap:conformity-rsvp}
The harmonic conformity relation $\Gamma_n(t)$ is
reconstructible as a constraint propagation indicator:
$\Gamma_n(t) = 1$ when harmonic $n$ is evolving within
the admissible trajectory basin of the fundamental,
and $\Gamma_n(t) = 0$ when it is not. This formulation
connects the \HCFone{} compression result to the RSVP
principle of constraint-based representation efficiency,
preserving the engineering result while making its
theoretical assumptions explicit.
\end{auditprop}

\section{The Spectral Fracture Hypothesis Revisited}
\label{sec:spectral-fracture-rsvp}

Chapter~5 noted that the spectral fracture concept ---
the claim that lossy compression disrupts temporal
relationships rather than merely reducing spectral
content --- was phenomenologically real but definitionally
risky. The RSVP translation provides a more careful
formulation.

In constraint-propagation terms, what block-based
transform coding disrupts is not the within-block
trajectory structure (which the codec preserves
reasonably well within the psychoacoustic model) but
the across-block constraint relationships. The
conformity relation between harmonics is computed
from amplitude envelope dynamics $v_n(t) = dA_n/dt / A_n$.
At block boundaries, the decoded waveform's amplitude
envelopes are determined by the spectral coefficients
within each block independently; the relationship
between the envelope dynamics at the end of one block
and the beginning of the next is not explicitly preserved.

This is a specific, testable claim: the \HCFone{}
conformity rate should be lower in compressed-and-decoded
audio at block boundaries than in the original, and
the drop should be proportional to the aggressiveness
of the compression. This prediction is at \LevelI{}
and does not require the authenticity ontology or the
consciousness hypothesis to be stated. It is the
recoverable core of the spectral fracture claim.

\begin{epistemicstatus}{Operationally Grounded}
The claim that lossy compression reduces conformity
rates at block boundaries, as measured by the \HCFone{}
conformity relation, is a \LevelI{} prediction derivable
from the block-based architecture of standard codecs
and the definition of the conformity relation. It is
testable with existing tools and datasets. This is
the constraint-preserving reformulation of the spectral
fracture claim: it preserves the engineering content
while discarding the authenticity ontology.
\end{epistemicstatus}

\section{What the Signal-Level Reconstruction Achieves}
\label{sec:signal-reconstruction-achievement}

The reconstruction of the signal-level components
achieves three things.

First, it separates the engineering results from the
philosophical commitments they were bundled with in
the original Phoenix formulation. The jitter metric,
the conformity relation, and the compression predictions
are all present in $\ThetaStar$, unchanged in their
engineering content. The consciousness claims, the
authenticity ontology, and the universality ambitions
are absent --- not because they are false, but because
they are not required to state the engineering results.

Second, it makes the research directions generated
by the engineering results explicit. The proportionality
$\SalFunc \propto 1/\AccEnt$ is not a derivation; it
is a hypothesis that could be tested by specifying
the state space for audio signals and computing both
quantities. The constraint propagation interpretation
of the conformity relation suggests a connection to
the RSVP dynamics that could be explored formally.
These are genuine research directions generated by
the reconstruction.

Third, it provides a vocabulary in which the boundary
conditions of the engineering results can be stated
clearly. The Marine metric is a low-entropy trajectory
detector for audio signals in time-domain representation;
it is not a general-purpose meaning detector, a
consciousness metric, or an authenticity test. The
RSVP translation makes this boundary explicit in its
terminology rather than leaving it implicit in the
reader's interpretation.
